% =====================================================================
%  Banco complementar --- Potencia Eletrica e Energia  (REVISADO)
%  Substitui a versao gerada por template (7 clones em N2, 8 em N3 e
%  8 questoes conceituais marcadas como N4).
%  O conteudo de ENERGIA esta no cap02, que ja cobre: rendimento
%  100 W/95 W, tarifa media a partir de conta de R\$ 80 / 200 kWh,
%  ferro eletrico e outros aparelhos por 2 h, e o consumo diario
%  residencial (lampada, chuveiro etc.). O cap05 cobre P=RI^2, P=VI
%  e maxima transferencia. O banco de circuitos serie revisado ja usa
%  um calculo de custo mensal (300 W, 6 h/dia, 30 dias).
%  Este banco explora: potencia MEDIA x valor eficaz, regime pulsado
%  e constante termica, rendimento em cascata, tensao errada de
%  alimentacao, energia em transitorio exponencial, payback e o
%  dimensionamento economico de condutor.
%  Distribuicao mantida: 5 N1 + 7 N2 + 8 N3 + 8 N4 = 28
% =====================================================================

\subsubsection*{Banco complementar --- Potência elétrica e energia}

\begin{atencao}[Organização do banco]
As três formas de $P$ ($UI$, $RI^{2}$ e $U^{2}/R$) são equivalentes apenas
quando vale a lei de Ohm; escolha a que usa os dados disponíveis sem passos
intermediários. Para \emph{energia}, lembre que $\SI{1}{\kilo\watt\hour}
=\SI{3.6}{\mega\joule}$. O N3 trata de regime pulsado, rendimento em cascata e
tensão incorreta de alimentação; o N4 exige integração, demonstração de que a
potência média usa o valor \emph{eficaz}, e análises de retorno econômico.
\end{atencao}

\blocofixacao

\begin{questao}[1]{}
Um resistor de $\SI{8}{\ohm}$ é percorrido por $\SI{2.5}{\ampere}$. Determine a
potência dissipada.
\resposta{$P=\SI{50}{\watt}$}
\resolucao{
\[
P=RI^{2}=8(2{,}5)^{2}=8(6{,}25)=\SI{50}{\watt}.
\]
\textbf{Escolha da fórmula:} conhecendo $R$ e $I$, use $RI^{2}$ diretamente.
Calcular $U=RI=\SI{20}{\volt}$ e depois $P=UI$ dá o mesmo resultado, mas
acrescenta um passo e uma oportunidade de erro.}
\end{questao}

\begin{questao}[1]{}
Um resistor de $\SI{44}{\ohm}$ é submetido a $\SI{220}{\volt}$. Determine a
potência.
\resposta{$P=\SI{1100}{\watt}$}
\resolucao{
\[
P=\frac{U^{2}}{R}=\frac{220^{2}}{44}=\frac{48400}{44}=\SI{1100}{\watt}.
\]
Note a dependência \emph{quadrática} na tensão: uma queda de $10\%$ na rede
($\SI{198}{\volt}$) reduziria a potência para
$\SI{891}{\watt}$ --- $19\%$ a menos. É por isso que chuveiros "esquentam
menos" em horários de pico.}
\end{questao}

\begin{questao}[1]{}
Converta $\SI{250}{\watt\hour}$ para joules.
\resposta{$\SI{900}{\kilo\joule}$}
\resolucao{
\[
E=Pt=250\,\si{\watt}\times3600\,\si{\second}=\SI{900000}{\joule}
=\SI{900}{\kilo\joule}.
\]
\textbf{Fator de conversão útil:}
$\SI{1}{\kilo\watt\hour}=\SI{3.6}{\mega\joule}$. O joule é pequeno demais para
contas de energia elétrica --- daí o uso do quilowatt-hora, que é uma unidade
de energia, não de potência. Confundir os dois é o erro mais comum do tópico.}
\end{questao}

\begin{questao}[1]{}
Conhecendo apenas a tensão sobre um resistor e o valor de sua resistência, a
expressão adequada para a potência é:
\begin{alternativas}[2]
\item $P=UI$
\item $P=RI^{2}$
\item $P=U^{2}/R$
\item $P=U/R$
\end{alternativas}
\resposta{(c)}
\resolucao{As alternativas (a) e (b) exigiriam calcular $I$ antes. A (d) está
dimensionalmente errada: $\si{\volt}/\si{\ohm}$ é ampère, não watt.\\
\textbf{Verificação dimensional rápida:} $U^{2}/R$ dá
$\si{\volt\squared}/\si{\ohm}=\si{\volt}\cdot\si{\ampere}=\si{\watt}$
\checkmark. Sempre que houver dúvida entre fórmulas parecidas, confira a
dimensão --- ela elimina a maioria das alternativas erradas.}
\end{questao}

\begin{questao}[1]{}
Um dispositivo de $\SI{60}{\watt}$ opera por $\SI{5}{\minute}$. Determine a
energia consumida em joules.
\resposta{$E=\SI{18}{\kilo\joule}$}
\resolucao{Convertendo o tempo para segundos:
\[
E=Pt=60\times(5\times60)=60(300)=\SI{18000}{\joule}=\SI{18}{\kilo\joule}.
\]
\textbf{Cuidado com as unidades:} $P$ em watts exige $t$ em \emph{segundos}
para obter joules. Se o tempo ficar em minutos, o resultado sai
$60\times$ menor --- erro silencioso, porque o número parece plausível.}
\end{questao}

\blocoaplicacao

\begin{questao}[2]{}
Um chuveiro de $\SI{5500}{\watt}$ opera em $\SI{220}{\volt}$. Determine a
corrente e a resistência.
\resposta{$I=\SI{25}{\ampere}$; $R=\SI{8.8}{\ohm}$}
\resolucao{
\[
I=\frac{P}{U}=\frac{5500}{220}=\SI{25}{\ampere};
\qquad
R=\frac{U^{2}}{P}=\frac{48400}{5500}=\SI{8.8}{\ohm}.
\]
\textbf{Conferência:} $R=U/I=220/25=\SI{8.8}{\ohm}$ \checkmark.\\
\textbf{Implicação de instalação:} $\SI{25}{\ampere}$ é uma corrente elevada.
Ela exige circuito exclusivo, condutor de seção adequada e disjuntor de
$\SI{32}{\ampere}$ --- e é por isso que chuveiro nunca compartilha circuito com
tomadas.}
\end{questao}

\begin{questao}[2]{}
Projete a resistência de um aquecedor de $\SI{1200}{\watt}$ para
$\SI{220}{\volt}$ e determine a corrente nominal e o disjuntor adequado.
\resposta{$R=\SI{40.3}{\ohm}$; $I=\SI{5.45}{\ampere}$; disjuntor de
$\SI{10}{\ampere}$}
\resolucao{
\[
R=\frac{U^{2}}{P}=\frac{48400}{1200}=\SI{40.3}{\ohm};
\qquad
I=\frac{1200}{220}=\SI{5.45}{\ampere}.
\]
\textbf{Disjuntor:} o valor comercial imediatamente acima é
$\SI{10}{\ampere}$, que oferece folga de $83\%$ --- adequado, pois aquecedores
resistivos não têm corrente de partida elevada.\\
\textbf{Nota sobre o material:} $\SI{40.3}{\ohm}$ deve ser obtido com fio de
liga resistiva (níquel-cromo) de comprimento e seção calculados. E como a
resistência desses materiais varia pouco com a temperatura, a potência
permanece próxima do nominal em regime.}
\end{questao}

\begin{questao}[2]{}
Duas lâmpadas incandescentes de $\SI{60}{\watt}$ são fabricadas para
$\SI{127}{\volt}$ e para $\SI{220}{\volt}$. Determine a resistência de cada
uma.
\resposta{$\SI{268.8}{\ohm}$ e $\SI{806.7}{\ohm}$}
\resolucao{
\[
R_{127}=\frac{127^{2}}{60}=\SI{268.8}{\ohm};
\qquad
R_{220}=\frac{220^{2}}{60}=\SI{806.7}{\ohm}.
\]
A lâmpada de tensão maior tem resistência \textbf{três vezes} maior --- razão
$(220/127)^{2}=3{,}0$.\\
\textbf{Consequência construtiva:} para a mesma potência, o filamento de
$\SI{220}{\volt}$ é mais longo e mais fino. Ele é, portanto, mais frágil
mecanicamente --- fato que se observa na prática, com lâmpadas de tensão maior
tendendo a durar menos sob vibração.}
\end{questao}

\begin{questao}[2]{}
Um aparelho opera em dois modos: $\SI{1500}{\watt}$ durante
$\SI{20}{\minute}$ e $\SI{600}{\watt}$ durante $\SI{40}{\minute}$, todos os
dias. Determine o consumo mensal ($\SI{30}{\day}$) e o custo a
$\mathrm{R\$}\,0{,}95$ por $\si{\kilo\watt\hour}$.
\resposta{$\SI{27}{\kilo\watt\hour}$; $\mathrm{R\$}\,25{,}65$}
\resolucao{Energia diária, com os tempos em horas:
\[
E=1{,}5\left(\tfrac13\right)+0{,}6\left(\tfrac23\right)
=0{,}50+0{,}40=\SI{0.90}{\kilo\watt\hour}.
\]
\[
E_{\text{mês}}=0{,}90(30)=\SI{27}{\kilo\watt\hour};
\qquad
\text{custo}=27(0{,}95)=\mathrm{R\$}\,25{,}65.
\]
\textbf{Observação:} o modo de $\SI{600}{\watt}$, embora tenha menos da metade
da potência, responde por $44\%$ do consumo --- porque funciona o dobro do
tempo. Em energia, \textbf{potência sozinha não decide}; o produto pelo tempo é
que conta.}
\end{questao}

\begin{questao}[2]{}
Um resistor é submetido a pulsos de $\SI{50}{\watt}$ com ciclo de trabalho de
$\SI{10}{\percent}$. Determine a potência média.
\resposta{$P_{\text{média}}=\SI{5}{\watt}$}
\resolucao{
\[
P_{\text{média}}=P_{pico}\times D=50(0{,}10)=\SI{5}{\watt}.
\]
\textbf{Mas cuidado:} isso \emph{não} significa que um resistor de
$\SI{5}{\watt}$ sirva. Se a duração do pulso for comparável ou maior que a
constante térmica do componente, ele atinge temperaturas próximas às de
$\SI{50}{\watt}$ contínuos e pode falhar.\\
A questão de quando a potência média basta é retomada quantitativamente no bloco
de aprofundamento.}
\end{questao}

\begin{questao}[2]{}
Uma fonte com rendimento de $\SI{90}{\percent}$ alimenta um conversor de
$\SI{85}{\percent}$. Determine o rendimento global e a potência de entrada
necessária para entregar $\SI{100}{\watt}$ úteis.
\resposta{$\eta=76{,}5\%$; $P_{entrada}=\SI{130.7}{\watt}$}
\resolucao{Rendimentos em cascata \textbf{multiplicam-se}:
\[
\eta=0{,}90\times0{,}85=0{,}765=76{,}5\%.
\]
\[
P_{entrada}=\frac{100}{0{,}765}=\SI{130.7}{\watt}.
\]
\textbf{Erro comum:} \emph{somar} ou fazer a média dos rendimentos, o que daria
$87{,}5\%$ --- valor otimista demais. Cada estágio desperdiça uma fração do que
recebe, e as perdas se compõem multiplicativamente.\\
\textbf{Perdas:} $\SI{13.1}{\watt}$ na fonte e $\SI{17.6}{\watt}$ no
conversor, totalizando $\SI{30.7}{\watt}$ dissipados.}
\end{questao}

\begin{questao}[2]{}
Uma carga alterna entre $\SI{500}{\watt}$ por $\SI{5}{\minute}$ e
$\SI{100}{\watt}$ por $\SI{15}{\minute}$, em ciclos de
$\SI{20}{\minute}$. Determine a potência média e a energia por ciclo.
\resposta{$P_{\text{média}}=\SI{200}{\watt}$; $E=\SI{66.7}{\watt\hour}$}
\resolucao{A potência média é a \textbf{média ponderada pelo tempo}:
\[
P_{\text{média}}=\frac{500(5)+100(15)}{20}=\frac{2500+1500}{20}=\SI{200}{\watt}.
\]
\[
E_{ciclo}=P_{\text{média}}\times t_{ciclo}=200\times\frac{20}{60}
=\SI{66.7}{\watt\hour}.
\]
\textbf{Atenção:} a média \emph{aritmética} das potências ($300\,\si{\watt}$)
estaria errada --- ela ignoraria que o patamar baixo dura três vezes mais.
Média de potência é sempre ponderada pelo tempo.}
\end{questao}

\blocoaprofundamento

\begin{questao}[3]{}
Um resistor de $\SI{10}{\ohm}$ conduz $\SI{4}{\ampere}$ durante
$\SI{25}{\percent}$ do tempo e nada nos $\SI{75}{\percent}$ restantes.
\begin{itensq}
\item Determine a corrente média e a potência calculada a partir dela.
\item Determine a potência média verdadeira.
\item Explique a discrepância.
\end{itensq}
\resposta{(a) $\SI{1}{\ampere}$ e $\SI{10}{\watt}$ (errado); (b)
$\SI{40}{\watt}$; (c) a potência depende de $\langle i^{2}\rangle$}
\resolucao{(a) $I_{\text{média}}=0{,}25(4)+0{,}75(0)=\SI{1}{\ampere}$, e
$RI_{\text{média}}^{2}=10(1)=\SI{10}{\watt}$.\\
(b) A potência instantânea vale $\SI{160}{\watt}$ durante $25\%$ do tempo e
zero no resto:
\[
P_{\text{média}}=0{,}25(160)+0{,}75(0)=\SI{40}{\watt}.
\]
\textbf{Quatro vezes} o valor obtido em (a).\\
(c) A potência é \emph{quadrática} na corrente, e a média de um quadrado não é
o quadrado da média:
\[
P_{\text{média}}=R\left\langle i^{2}\right\rangle
\ne R\left\langle i\right\rangle^{2}.
\]
Define-se então o \textbf{valor eficaz}
$I_{rms}=\sqrt{\langle i^{2}\rangle}=\sqrt{0{,}25(16)}=\SI{2}{\ampere}$, com o
qual $P=RI_{rms}^{2}=\SI{40}{\watt}$ \checkmark.\\
\textbf{Definição operacional do valor eficaz:} é a corrente contínua que
produziria o \emph{mesmo aquecimento}. Aqui, $\SI{2}{\ampere}$ contínuos
aqueceriam tanto quanto o regime pulsado --- e é por isso que multímetros de
qualidade medem "true RMS".}
\end{questao}

\begin{questao}[3]{}
Um resistor com potência nominal de $\SI{10}{\watt}$ recebe pulsos de
$\SI{50}{\watt}$ com ciclo de trabalho de $\SI{10}{\percent}$.
\begin{itensq}
\item Determine a potência média e compare com o nominal.
\item Estabeleça o critério para decidir se o componente sobrevive.
\item Analise dois casos: pulsos de $\SI{1}{\milli\second}$ e de
      $\SI{10}{\second}$.
\end{itensq}
\resposta{$P_{\text{média}}=\SI{5}{\watt}$; a decisão depende da constante térmica}
\resolucao{(a) $P_{\text{média}}=\SI{5}{\watt}$, metade do nominal --- aparentemente
seguro.\\
(b) \textbf{O critério é a comparação entre a duração do pulso e a constante
térmica $\tau_t$ do componente} (tipicamente de segundos a dezenas de segundos
para resistores de fio; menos para os de filme).
\begin{itemize}
\item Se $t_{pulso}\ll\tau_t$, a massa térmica \emph{integra} os pulsos: a
temperatura acompanha a potência \textbf{média}, e o critério de
$\SI{5}{\watt}$ vale.
\item Se $t_{pulso}\gg\tau_t$, o resistor atinge o equilíbrio térmico
\emph{durante} cada pulso: a temperatura corresponde a $\SI{50}{\watt}$
contínuos, e o componente queima.
\end{itemize}
(c) \textbf{Pulsos de $\SI{1}{\milli\second}$:} muito menores que qualquer
$\tau_t$ realista --- o resistor "enxerga" $\SI{5}{\watt}$. \textbf{Seguro.}\\
\textbf{Pulsos de $\SI{10}{\second}$:} comparáveis ou maiores que $\tau_t$ --- o
resistor esquenta como se recebesse $\SI{50}{\watt}$ por dez segundos.
\textbf{Falha provável.}\\
\textbf{Consequência prática:} fabricantes publicam curvas de \emph{potência de
pulso} justamente porque o valor nominal contínuo não responde a essa pergunta.
Nunca dimensione regime pulsado apenas pela média.}
\end{questao}

\begin{questao}[3]{}
Um sistema tem uma fonte de $\SI{92}{\percent}$ em cascata com um conversor de
$\SI{85}{\percent}$.
\begin{itensq}
\item Determine o rendimento global.
\item Compare o ganho de melhorar o estágio pior (de $85$ para
      $\SI{92}{\percent}$) com o de melhorar o melhor (de $92$ para
      $\SI{97}{\percent}$).
\item Formule a regra geral.
\end{itensq}
\resposta{(a) $78{,}2\%$; (b) $84{,}6\%$ contra $82{,}5\%$ --- melhorar o pior
rende mais}
\resolucao{(a) $\eta=0{,}92(0{,}85)=0{,}782=78{,}2\%$.\\
(b) \textbf{Melhorando o pior:} $0{,}92(0{,}92)=84{,}6\%$ --- ganho de
$6{,}4$ pontos.\\
\textbf{Melhorando o melhor:} $0{,}97(0{,}85)=82{,}5\%$ --- ganho de
$4{,}3$ pontos.\\
E note que o segundo caso exigiu uma melhoria maior em pontos absolutos
($5$ contra $7$) para render \emph{menos}.\\
(c) \textbf{Regra geral.} Como $\eta=\prod\eta_k$, o ganho relativo global é a
soma dos ganhos relativos:
\[
\frac{\Delta\eta}{\eta}=\sum_k\frac{\Delta\eta_k}{\eta_k}.
\]
Um mesmo acréscimo \emph{absoluto} rende mais no estágio de menor $\eta_k$,
pois ali representa maior acréscimo relativo. \textbf{Ataque sempre o elo mais
fraco} --- e é também nele que a margem de melhoria costuma ser maior e mais
barata.}
\end{questao}

\begin{questao}[3]{}
Um aparelho consome $\SI{4}{\watt}$ em espera, permanentemente.
\begin{itensq}
\item Determine o consumo anual e o custo a
      $\mathrm{R\$}\,0{,}95/\si{\kilo\watt\hour}$.
\item Estime o impacto de dez aparelhos semelhantes.
\item Proponha uma decisão de engenharia.
\end{itensq}
\resposta{$\SI{35}{\kilo\watt\hour}$ e $\mathrm{R\$}\,33{,}29$ por aparelho;
$\mathrm{R\$}\,333$ para dez}
\resolucao{(a) Um ano tem $\SI{8760}{\hour}$:
\[
E=4\times8760=\SI{35040}{\watt\hour}=\SI{35.04}{\kilo\watt\hour};
\]
\[
\text{custo}=35{,}04(0{,}95)=\mathrm{R\$}\,33{,}29.
\]
(b) Dez aparelhos: $\SI{350}{\kilo\watt\hour}$ e cerca de
$\mathrm{R\$}\,333$ por ano --- comparável ao consumo anual de uma geladeira
eficiente.\\
(c) \textbf{Decisões possíveis, com seus custos.}
\begin{itemize}
\item \textbf{Filtro com chave} para desligar grupos de aparelhos: custo de
dezenas de reais, retorno em poucos meses. Melhor relação custo-benefício.
\item \textbf{Redesenho do produto:} normas modernas exigem espera abaixo de
$\SI{0.5}{\watt}$, o que reduziria o custo a menos de
$\mathrm{R\$}\,4{,}20$ por ano.
\item \textbf{Não fazer nada:} defensável para um aparelho isolado, insustentável
em escala --- o consumo agregado de espera responde por vários por cento do
consumo residencial nacional.
\end{itemize}
\textbf{O ponto de engenharia:} $\SI{4}{\watt}$ parecem irrelevantes porque a
potência é pequena. Mas energia é potência \emph{vezes tempo}, e
$\SI{8760}{\hour}$ é um multiplicador enorme.}
\end{questao}

\begin{questao}[3]{}
Um aparelho projetado para $\SI{220}{\volt}$ e $\SI{1000}{\watt}$ é ligado por
engano em $\SI{127}{\volt}$.
\begin{itensq}
\item Determine a potência resultante, supondo resistência constante.
\item Analise o caso inverso: aparelho de $\SI{127}{\volt}$ em
      $\SI{220}{\volt}$.
\end{itensq}
\resposta{(a) $\SI{333}{\watt}$ ($33\%$); (b) potência \emph{triplicada} ---
destruição}
\resolucao{(a) $R=\dfrac{220^{2}}{1000}=\SI{48.4}{\ohm}$, e em
$\SI{127}{\volt}$:
\[
P=\frac{127^{2}}{48{,}4}=\SI{333}{\watt}
=\left(\frac{127}{220}\right)^{2}\times1000=33{,}3\%.
\]
\textbf{Consequência:} o aparelho funciona, mas com um terço da potência --- um
chuveiro esquenta pouco, um aquecedor não aquece. Situação incômoda, porém
\emph{segura}.\\
(b) $P'=\left(\dfrac{220}{127}\right)^{2}P=3{,}0\,P$ --- a potência
\textbf{triplica}, e a corrente quase dobra.\\
\textbf{Resultado:} superaquecimento imediato, fumaça e, muitas vezes, incêndio.
O disjuntor pode nem atuar, pois a corrente ($\SI{13.6}{\ampere}$ para um
aparelho de $\SI{1000}{\watt}$/$\SI{127}{\volt}$) fica abaixo do valor de
disparo --- a proteção existe para o \emph{circuito}, não para o aparelho.\\
\textbf{Assimetria fundamental:} subtensão é inconveniente; sobretensão é
destrutiva. Por isso a etiqueta de tensão merece conferência antes de ligar
qualquer aparelho resistivo.}
\end{questao}

\begin{questao}[3]{}
Um resistor de $\SI{4}{\ohm}$ é percorrido por
$i(t)=\SI{5}{\ampere}\cdot e^{-t/\tau}$, com $\tau=\SI{20}{\milli\second}$.
\begin{itensq}
\item Determine a potência inicial.
\item Determine a energia total dissipada.
\item Determine o instante em que metade da energia já foi dissipada.
\end{itensq}
\resposta{(a) $\SI{100}{\watt}$; (b) $\SI{1}{\joule}$; (c)
$t=\SI{6.93}{\milli\second}$}
\resolucao{(a) $P_0=RI_0^{2}=4(25)=\SI{100}{\watt}$.\\
(b) A potência decai como $e^{-2t/\tau}$:
\[
E=\int_0^{\infty}RI_0^{2}e^{-2t/\tau}\,\mathrm{d}t
=RI_0^{2}\cdot\frac{\tau}{2}
=100\times\frac{0{,}020}{2}=\SI{1}{\joule}.
\]
\textbf{Regra útil:} a energia equivale à potência inicial mantida por
$\tau/2$. A "constante de tempo da potência" é metade da constante de tempo da
corrente, porque a potência vai com o quadrado.\\
(c) Metade da energia:
\[
\int_0^{t}e^{-2s/\tau}\mathrm{d}s=\frac12\int_0^{\infty}
\;\Longrightarrow\;
1-e^{-2t/\tau}=\frac12
\;\Longrightarrow\;
t=\frac{\tau\ln2}{2}=\SI{6.93}{\milli\second}.
\]
\textbf{Interpretação:} metade de toda a energia é liberada em apenas
$0{,}35\,\tau$ --- o transitório é fortemente concentrado no início. É o pico
inicial, e não a duração, que dimensiona o componente.}
\end{questao}

\begin{questao}[3]{}
Compare uma lâmpada incandescente de $\SI{60}{\watt}$ com uma de LED de
$\SI{9}{\watt}$ e mesmo fluxo luminoso, usadas $\SI{5}{\hour}$ por dia.
\begin{itensq}
\item Determine o consumo e o custo anual de cada uma, a
      $\mathrm{R\$}\,0{,}95/\si{\kilo\watt\hour}$.
\item Determine o tempo de retorno se a lâmpada de LED custar
      $\mathrm{R\$}\,15{,}00$ a mais.
\end{itensq}
\resposta{(a) $\mathrm{R\$}\,104{,}02$ contra $\mathrm{R\$}\,15{,}60$;
(b) cerca de $2 meses$}
\resolucao{(a) Horas anuais: $5\times365=\SI{1825}{\hour}$.
\[
E_{inc}=0{,}060(1825)=\SI{109.5}{\kilo\watt\hour}
\Rightarrow \mathrm{R\$}\,104{,}02;
\]
\[
E_{LED}=0{,}009(1825)=\SI{16.4}{\kilo\watt\hour}
\Rightarrow \mathrm{R\$}\,15{,}60.
\]
Economia anual: $\mathrm{R\$}\,88{,}42$.\\
(b)
\[
t=\frac{15{,}00}{88{,}42}=0{,}17\ \text{ano}\approx2 meses.
\]
\textbf{Retorno excepcional} --- poucos investimentos em engenharia se pagam em
dois meses.\\
\textbf{E há mais:} a lâmpada de LED dura tipicamente $\SI{25000}{\hour}$
contra $\SI{1000}{\hour}$ da incandescente, evitando cerca de $24$ trocas ao
longo da vida. Somando o custo das lâmpadas substituídas e da mão de obra, a
vantagem cresce ainda mais.\\
\textbf{Nota física:} a incandescente converte cerca de $95\%$ da energia em
calor e apenas $5\%$ em luz --- ela é, essencialmente, um aquecedor que
ilumina de lado.}
\end{questao}

\begin{questao}[3]{}
Uma instalação alimenta uma carga de $\SI{10}{\kilo\watt}$ por um cabo de
resistência total $\SI{0.2}{\ohm}$, em $\SI{220}{\volt}$.
\begin{itensq}
\item Determine a corrente e a perda no cabo.
\item Determine a fração da energia perdida.
\item Repita para $\SI{440}{\volt}$ e comente.
\end{itensq}
\resposta{(a) $\SI{45.5}{\ampere}$, $\SI{414}{\watt}$; (b) $4{,}1\%$;
(c) $\SI{103}{\watt}$, $1{,}0\%$}
\resolucao{(a) $I=\dfrac{10000}{220}=\SI{45.5}{\ampere}$ e
\[
P_{cabo}=RI^{2}=0{,}2(45{,}5)^{2}=\SI{414}{\watt}.
\]
(b) $414/10000=4{,}1\%$ da potência útil vira calor no cabo.\\
(c) Em $\SI{440}{\volt}$: $I=\SI{22.7}{\ampere}$ e
$P_{cabo}=0{,}2(22{,}7)^{2}=\SI{103}{\watt}$, ou $1{,}0\%$.\\
\textbf{Dobrar a tensão reduziu a perda a um quarto.} A razão é direta: para a
mesma potência, $I\propto1/U$, e a perda vai com $I^{2}$ --- logo
$P_{perda}\propto1/U^{2}$.\\
\textbf{É esta a razão de existir a alta tensão na transmissão.} Linhas de
$\SI{500}{\kilo\volt}$ transportam gigawatts com perdas de poucos por cento; nas
mesmas linhas em baixa tensão, a energia não chegaria ao destino. O preço é o
isolamento e os transformadores nas duas pontas.}
\end{questao}

\blocodesafio

\begin{questao}[4]{}
\textbf{(Demonstração --- potência média e valor eficaz.)} Prove que a potência
média em um resistor depende do valor \emph{eficaz} da corrente.
\begin{itensq}
\item Deduza $P_{\text{média}}=R\,I_{rms}^{2}$ a partir da definição.
\item Mostre que $I_{rms}\ge I_{\text{média}}$, com igualdade apenas para corrente
      constante.
\item Aplique a uma onda quadrada de $\SI{0}{\ampere}$ e
      $\SI{4}{\ampere}$ com ciclo de $\SI{25}{\percent}$.
\end{itensq}
\resposta{$I_{rms}=\SI{2}{\ampere}$ contra $I_{\text{média}}=\SI{1}{\ampere}$;
$P=R I_{rms}^{2}$}
\resolucao{(a) Por definição, a potência média em um período $T$ é
\[
P_{\text{média}}=\frac{1}{T}\int_0^{T}R\,i^{2}(t)\,\mathrm{d}t
=R\left[\frac{1}{T}\int_0^{T}i^{2}\,\mathrm{d}t\right]
=R\,I_{rms}^{2},
\]
onde a definição de valor eficaz é justamente
$I_{rms}=\sqrt{\frac1T\int_0^{T}i^{2}\mathrm{d}t}$ --- a
\emph{raiz} da \emph{média} do \emph{quadrado}.\\
(b) Pela desigualdade de Cauchy--Schwarz (ou pelo fato de a variância ser não
negativa):
\[
\left\langle i^{2}\right\rangle-\left\langle i\right\rangle^{2}
=\operatorname{Var}(i)\ge0
\;\Longrightarrow\;
I_{rms}^{2}\ge I_{\text{média}}^{2}.
\]
A igualdade vale se e somente se a variância é nula, isto é, se $i(t)$ é
\textbf{constante}.\\
\textbf{Leitura:} qualquer flutuação da corrente aumenta o aquecimento em
relação ao que a média sugeriria. Corrente ondulada aquece mais que corrente
contínua de mesmo valor médio --- resultado com consequências diretas em
retificadores e conversores.\\
(c) $\langle i\rangle=0{,}25(4)=\SI{1}{\ampere}$;
$\langle i^{2}\rangle=0{,}25(16)=\SI{4}{\ampere\squared}$, logo
$I_{rms}=\SI{2}{\ampere}$.\\
Com $R=\SI{10}{\ohm}$: $P=10(4)=\SI{40}{\watt}$, e não os
$\SI{10}{\watt}$ que a corrente média sugeriria. Fator $4$ --- exatamente
$(I_{rms}/I_{\text{média}})^{2}$.}
\end{questao}

\begin{questao}[4]{}
\textbf{(Integração --- energia em transitório.)} Um resistor $R$ conduz
$i(t)=I_0e^{-t/\tau}$.
\begin{itensq}
\item Deduza a energia total dissipada.
\item Determine a "potência média equivalente" ao longo de $5\tau$.
\item Compare com a energia de um pulso retangular de mesma amplitude e
      duração $\tau$.
\end{itensq}
\resposta{$E=\dfrac{RI_0^{2}\tau}{2}$; equivale à potência inicial mantida por
$\tau/2$}
\resolucao{(a)
\[
E=\int_0^{\infty}R\,I_0^{2}e^{-2t/\tau}\,\mathrm{d}t
=RI_0^{2}\left[-\frac{\tau}{2}e^{-2t/\tau}\right]_0^{\infty}
=\frac{RI_0^{2}\tau}{2}.
\]
\textbf{Observe o fator $\tau/2$, e não $\tau$:} a potência decai com constante
$\tau/2$ porque é proporcional a $i^{2}$, e o quadrado de uma exponencial tem
metade da constante de tempo.\\
(b) Distribuindo $E$ por $5\tau$:
\[
P_{\text{média}}=\frac{RI_0^{2}\tau/2}{5\tau}=\frac{RI_0^{2}}{10}=\frac{P_0}{10}.
\]
Apenas um décimo da potência inicial --- confirmando que o transitório é curto
diante da janela considerada.\\
(c) \textbf{Pulso retangular} de amplitude $I_0$ e duração $\tau$:
$E_{ret}=RI_0^{2}\tau$, exatamente o \textbf{dobro}.\\
\textbf{Consequência de dimensionamento:} substituir o decaimento real por um
pulso retangular equivalente de duração $\tau$ é uma aproximação
\emph{conservadora} --- ela superestima a energia por um fator $2$. Útil como
estimativa rápida e segura, desde que se saiba que há folga embutida.}
\end{questao}

\begin{questao}[4]{}
\textbf{(Projeto de proteção.)} Um resistor de $\SI{100}{\ohm}$ e
$\frac14\,\si{\watt}$ deve ser protegido contra sobrecarga.
\begin{itensq}
\item Determine a tensão e a corrente máximas admissíveis.
\item Projete um resistor limitador em série, para uma fonte de
      $\SI{24}{\volt}$.
\item Verifique a potência do limitador.
\end{itensq}
\resposta{$U_{\max}=\SI{5}{\volt}$, $I_{\max}=\SI{50}{\milli\ampere}$;
limitador de $\SI{380}{\ohm}$}
\resolucao{(a)
\[
U_{\max}=\sqrt{PR}=\sqrt{0{,}25(100)}=\SI{5}{\volt};
\qquad
I_{\max}=\sqrt{\frac{P}{R}}=\sqrt{0{,}0025}=\SI{50}{\milli\ampere}.
\]
(b) Com $\SI{24}{\volt}$ e a exigência $I\le\SI{50}{\milli\ampere}$:
\[
R_{total}\ge\frac{24}{0{,}050}=\SI{480}{\ohm}
\;\Longrightarrow\;
R_{lim}\ge480-100=\SI{380}{\ohm}.
\]
O valor E24 imediatamente superior é $\SI{390}{\ohm}$, que dá
$I=24/490=\SI{49}{\milli\ampere}$ \checkmark.\\
(c) $P_{lim}=390(0{,}049)^{2}=\SI{0.94}{\watt}$ --- especifique
$\SI{2}{\watt}$.\\
\textbf{Observação incômoda:} o limitador dissipa \emph{quatro vezes} mais que o
resistor protegido, e o conjunto consome $\SI{1.18}{\watt}$ para entregar
$\SI{0.24}{\watt}$ ao componente de interesse.\\
\textbf{Quando isso se justifica:} apenas se o resistor de
$\SI{100}{\ohm}$ tiver função específica (sensor, elemento de calibração) que o
torne insubstituível. Caso contrário, redimensionar o próprio resistor para uma
potência maior é mais simples e mais eficiente.}
\end{questao}

\begin{questao}[4]{}
\textbf{(Análise de retorno.)} Dois equipamentos entregam a mesma energia útil
de $\SI{1000}{\kilo\watt\hour}$ por ano: o modelo A tem rendimento de
$\SI{80}{\percent}$ e o B, de $\SI{92}{\percent}$.
\begin{itensq}
\item Determine a energia de entrada de cada um e a diferença.
\item Determine o tempo de retorno se B custar $\mathrm{R\$}\,800$ a mais, com
      tarifa de $\mathrm{R\$}\,0{,}95/\si{\kilo\watt\hour}$.
\item Discuta os fatores que a análise simples ignora.
\end{itensq}
\resposta{(a) $1250$ e $\SI{1087}{\kilo\watt\hour}$, diferença de
$\SI{163}{\kilo\watt\hour}$; (b) cerca de $\SI{5.2}{year}$}
\resolucao{(a)
\[
E_A=\frac{1000}{0{,}80}=\SI{1250}{\kilo\watt\hour};
\qquad
E_B=\frac{1000}{0{,}92}=\SI{1087}{\kilo\watt\hour};
\]
diferença de $\SI{163}{\kilo\watt\hour}$ por ano.\\
(b) Economia anual: $163(0{,}95)=\mathrm{R\$}\,154{,}89$.
\[
t=\frac{800}{154{,}89}=\SI{5.2}{year}.
\]
(c) \textbf{O que a conta simples ignora.}
\begin{itemize}
\item \textbf{Reajuste da tarifa:} se a energia subir acima da inflação, o
retorno acelera. Historicamente é o que ocorre.
\item \textbf{Custo do capital:} $\mathrm{R\$}\,800$ aplicados renderiam algo
--- o retorno real é mais longo que o nominal.
\item \textbf{Vida útil:} se o equipamento durar $\SI{10}{year}$, B gera lucro
nos últimos cinco. Se durar $\SI{4}{year}$, a troca não se paga.
\item \textbf{Custo do calor rejeitado:} A dissipa $\SI{250}{\kilo\watt\hour}$
contra $\SI{87}{\kilo\watt\hour}$ de B. Em ambiente climatizado, esse calor
extra precisa ser \emph{removido} --- o que pode dobrar a economia efetiva.
\item \textbf{Regime de uso:} rendimentos variam com a carga, e o valor de
catálogo costuma ser o de plena carga.
\end{itemize}
\textbf{Conclusão:} $\SI{5.2}{year}$ é aceitável para equipamento industrial de
vida longa e uso contínuo; é ruim para uso intermitente ou equipamento de
substituição rápida. \textbf{A resposta depende do horizonte, não apenas do
número.}}
\end{questao}

\begin{questao}[4]{}
\textbf{(Restrição orçamentária.)} Deseja-se limitar a
$\mathrm{R\$}\,90{,}00$ o custo mensal de um aquecedor de
$\SI{1.5}{\kilo\watt}$, com tarifa de
$\mathrm{R\$}\,1{,}00/\si{\kilo\watt\hour}$.
\begin{itensq}
\item Determine o tempo de uso admissível.
\item Determine a potência que permitiria uso de $\SI{4}{\hour}$ diárias dentro
      do mesmo orçamento.
\item Discuta qual das duas soluções é preferível.
\end{itensq}
\resposta{(a) $\SI{60}{\hour}$ por mês, ou $\SI{2}{\hour}$ por dia;
(b) $\SI{750}{\watt}$}
\resolucao{(a) O orçamento permite $\SI{90}{\kilo\watt\hour}$ por mês:
\[
t=\frac{90}{1{,}5}=\SI{60}{\hour\per month}
=\SI{2}{\hour\per day}.
\]
(b) Para $4\times30=\SI{120}{\hour}$ mensais dentro dos mesmos
$\SI{90}{\kilo\watt\hour}$:
\[
P=\frac{90}{120}=\SI{0.75}{\kilo\watt}=\SI{750}{\watt}.
\]
(c) \textbf{Depende do objetivo térmico.} As duas soluções entregam a
\emph{mesma energia} --- e portanto o mesmo calor total no mês.
\begin{itemize}
\item \textbf{$\SI{1.5}{\kilo\watt}$ por $\SI{2}{\hour}$:} aquece rápido e
atinge temperatura mais alta, mas o ambiente esfria nas outras
$\SI{22}{\hour}$.
\item \textbf{$\SI{750}{\watt}$ por $\SI{4}{\hour}$:} aquecimento mais suave e
prolongado, com temperatura média mais estável.
\end{itemize}
\textbf{Para conforto térmico, a segunda costuma ser superior} --- o ambiente
tem inércia térmica, e potência menor por mais tempo reduz as oscilações.\\
\textbf{E há uma terceira solução, melhor que ambas:} melhorar o
\emph{isolamento} do ambiente. Isso reduz a energia necessária, e não apenas a
redistribui --- atacando a causa em vez do sintoma.}
\end{questao}

\begin{questao}[4]{}
\textbf{(Otimização econômica de condutor.)} Um cabo alimenta uma carga de
$\SI{40}{\ampere}$ por $\SI{4000}{\hour}$ ao ano. O custo do cabo é
proporcional à seção, e a resistência é inversamente proporcional a ela.
\begin{itensq}
\item Escreva o custo total como função da seção.
\item Determine a condição de mínimo.
\item Interprete o resultado.
\end{itensq}
\resposta{No ótimo, o custo anualizado do cabo iguala o custo anual da energia
perdida}
\resolucao{(a) Sejam $S$ a seção, $k_1$ o custo anualizado do cabo por unidade
de seção e $R=\rho\ell/S$. O custo total anual é
\[
C(S)=\underbrace{k_1S}_{\text{investimento}}
+\underbrace{\frac{\rho\ell}{S}I^{2}\,t\,c_E}_{\text{energia perdida}}
=k_1S+\frac{k_2}{S},
\]
com $k_2=\rho\ell I^{2}t\,c_E$.\\
(b) Derivando e anulando:
\[
\frac{\mathrm{d}C}{\mathrm{d}S}=k_1-\frac{k_2}{S^{2}}=0
\;\Longrightarrow\;
S^{*}=\sqrt{\frac{k_2}{k_1}}.
\]
A segunda derivada, $2k_2/S^{3}>0$, confirma o mínimo.\\
(c) \textbf{Substituindo $S^{*}$ nas duas parcelas:}
\[
k_1S^{*}=\sqrt{k_1k_2}
\qquad\text{e}\qquad
\frac{k_2}{S^{*}}=\sqrt{k_1k_2}.
\]
As duas são \textbf{iguais} no ótimo --- resultado conhecido como \emph{lei de
Kelvin}: \textbf{a seção econômica é aquela em que o custo anualizado do
condutor iguala o custo anual das perdas}.\\
\textbf{Leituras práticas.}
\begin{itemize}
\item Cabos usados \emph{muitas horas} por ano ($t$ grande) justificam seções
maiores --- $S^{*}\propto\sqrt{t}$.
\item Correntes maiores também: $S^{*}\propto I$.
\item Energia mais cara empurra a seção para cima: $S^{*}\propto\sqrt{c_E}$.
\end{itemize}
\textbf{Importante:} a seção econômica é frequentemente \emph{maior} que a
mínima exigida pela norma (que garante apenas segurança térmica e queda de
tensão). Dimensionar pelo mínimo normativo é seguro, mas costuma ser
economicamente subótimo em instalações de uso intenso.}
\end{questao}

\begin{questao}[4]{}
\textbf{(Balanço energético de um chuveiro.)} Um chuveiro de
$\SI{5500}{\watt}$ aquece água de $\SI{20}{\celsius}$ a
$\SI{40}{\celsius}$.
\begin{itensq}
\item Determine a vazão máxima, supondo rendimento de $\SI{100}{\percent}$.
\item Determine o custo de um banho de $\SI{8}{\minute}$ a
      $\mathrm{R\$}\,0{,}95/\si{\kilo\watt\hour}$.
\item Discuta por que o rendimento é tão alto neste caso.
\end{itensq}
\dados{$c_{\text{água}}=\SI{4180}{\joule\per\kilogram\per\kelvin}$}
\resposta{(a) $\SI{0.0658}{\kilogram\per\second}\approx
\SI{3.9}{\liter\per\minute}$; (b) $\mathrm{R\$}\,0{,}70$}
\resolucao{(a) A potência necessária para aquecer uma vazão mássica $\dot m$
por $\Delta T=\SI{20}{\kelvin}$ é $P=\dot m\,c\,\Delta T$:
\[
\dot m=\frac{P}{c\,\Delta T}=\frac{5500}{4180(20)}
=\SI{0.0658}{\kilogram\per\second}\approx\SI{3.9}{\liter\per\minute}.
\]
Vazão modesta --- e é exatamente por isso que abrir demais o registro faz a água
esfriar: a potência é fixa, e mais vazão significa menor $\Delta T$.\\
(b)
\[
E=5{,}5\,\si{\kilo\watt}\times\frac{8}{60}\,\si{\hour}
=\SI{0.733}{\kilo\watt\hour}
\;\Longrightarrow\;
\text{custo}=0{,}733(0{,}95)=\mathrm{R\$}\,0{,}70.
\]
Um banho diário custa cerca de $\mathrm{R\$}\,21$ por mês --- e para uma
família de quatro pessoas, mais de $\mathrm{R\$}\,80$.\\
(c) \textbf{Por que o rendimento é quase $100\%$.} A resistência está
\emph{imersa} na água, e \textbf{toda} a energia dissipada por efeito Joule vira
calor exatamente onde ele é desejado. Não há conversão intermediária, nem calor
"desperdiçado" --- o que seria perda em outro contexto é aqui o produto.\\
\textbf{Contraste instrutivo:} uma lâmpada incandescente tem os mesmos
$\approx100\%$ de conversão em calor, mas rendimento \emph{luminoso} de $5\%$.
O rendimento só se define em relação ao \emph{objetivo}: a mesma física é
excelente para aquecer e péssima para iluminar.}
\end{questao}

\begin{questao}[4]{}
\textbf{(Diagnóstico por consumo.)} A conta de energia de uma residência subiu
de $\SI{180}{\kilo\watt\hour}$ para $\SI{320}{\kilo\watt\hour}$ mensais, sem
mudança de hábitos.
\begin{itensq}
\item Determine a potência contínua equivalente ao acréscimo.
\item Levante hipóteses compatíveis com esse valor.
\item Proponha um procedimento de investigação.
\end{itensq}
\resposta{Acréscimo de $\SI{140}{\kilo\watt\hour}$, equivalente a
$\SI{194}{\watt}$ contínuos}
\resolucao{(a) O mês tem cerca de $\SI{720}{\hour}$:
\[
P=\frac{140000}{720}=\SI{194}{\watt}\ \text{contínuos}.
\]
(b) \textbf{Hipóteses compatíveis com $\approx\SI{200}{\watt}$ permanentes.}
\begin{itemize}
\item \textbf{Fuga para o terra} por isolamento danificado --- especialmente em
resistência de chuveiro ou aquecedor de água. Uma fuga de
$\SI{0.88}{\ampere}$ em $\SI{220}{\volt}$ dá exatamente
$\SI{194}{\watt}$.
\item \textbf{Geladeira com vedação ruim ou gás insuficiente}, operando em
regime quase contínuo em vez de ciclado.
\item \textbf{Bomba ou motor} funcionando sem parar por falha de boia ou
pressostato.
\item \textbf{Ligação clandestina} a jusante do medidor.
\end{itemize}
(c) \textbf{Procedimento de investigação, em ordem.}
\begin{enumerate}
\item \textbf{Desligar o disjuntor geral} e verificar se o medidor continua
girando. Se sim, o problema está \emph{antes} do quadro --- ligação irregular ou
medidor defeituoso.
\item \textbf{Religar circuito por circuito}, observando o medidor a cada
etapa, para isolar o ramo responsável.
\item \textbf{Medir a corrente} de cada circuito suspeito com alicate
amperímetro, com todos os aparelhos desligados. Corrente não nula indica fuga.
\item \textbf{Testar o dispositivo DR}, se houver --- fuga acima de
$\SI{30}{\milli\ampere}$ deveria ter provocado desarme. A ausência de desarme
com fuga confirmada indica DR defeituoso, e é problema de segurança, não apenas
de conta.
\end{enumerate}
\textbf{Prioridade:} uma fuga de $\SI{0.88}{\ampere}$ é risco de choque e de
incêndio. O custo na conta é o sintoma menos grave.}
\end{questao}
